% TOP_PROBLEM: {"problemId": "problem.collatz-conjecture", "canonicalTitle": "Collatz conjecture", "releaseRank": 188, "primaryDomain": "number_theory_arithmetic_geometry", "primaryDomainLabel": "Number theory & arithmetic geometry", "displayStatus": "open", "releaseStatus": "open"}
% !TeX program = lualatex
% Definition and sources only; this file is not an LLM solution attempt.
\documentclass[11pt]{article}
\usepackage[margin=25mm]{geometry}
\usepackage{fontspec}
\setmainfont{Latin Modern Roman}
\usepackage{amsmath,amssymb,mathtools}
\usepackage{unicode-math}
\setmathfont{Latin Modern Math}
\usepackage{xurl,hyperref}
\hypersetup{colorlinks=true,urlcolor=blue}
\setcounter{secnumdepth}{0}
\setlength{\parindent}{0pt}
\setlength{\parskip}{5pt}
\setlength{\emergencystretch}{3em}
\begin{document}
\section*{\texorpdfstring{Collatz conjecture}{Collatz conjecture}}\label{188}
\subsection{Definitions and mathematical statement}
Define $T:{\mathbb{N}}_{>0}\to{\mathbb{N}}_{>0}$ by
\[
 T(n)=\begin{cases}n/2,&2\mid n,\\3n+1,&2\nmid n.\end{cases}
\]
The Collatz assertion is $\forall n\ge1\ \exists k\ge0$ such that $T^k(n)=1$, where $T^0$ is the identity. Reaching one enters the cycle $1,4,2,1$. The iteration here uses the unaccelerated map; no division by an additional power of two is included in an odd step. A negative resolution may be an orbit entering a different cycle or one that never repeats and never reaches one. Verifying any finite initial range does not decide the universal assertion.
\par\smallskip\textit{Formulation and concept references:} \href{https://arxiv.org/abs/1909.03562}{[S1]}.

\subsection{Short English statement}
Does repeatedly halving even integers and replacing odd integers by three times the number plus one always eventually reach one?
\subsection{Sources}
\begin{itemize}
\item[S1]Terence Tao, Forum of Mathematics, Pi 10 (2022), e12.
\url{https://arxiv.org/abs/1909.03562}.
\textit{Formulation source.}
\end{itemize}
\end{document}
